\section{Examen General de Egreso (EGEL)}
Presentas el EGEL Plus del CENEVAL en sede UAEM y obtienes Desempeño Satisfactorio o Sobresaliente. (Art. 29, REP).

\subsection{Requisitos de la opción}
\begin{itemize}
  \item Tramitar la solicitud en el Departamento de Evaluación Profesional antes de presentar el examen.
  \item Presentar el EGEL en sede UAEM.
  \item Obtener testimonio de Desempeño Satisfactorio o Sobresaliente.
  \item Si el resultado es «Aún no satisfactorio» cuenta como evaluación aplazada; puedes pedir una segunda y última vez después de seis meses.
\end{itemize}
\subsection{Elegir la opción e iniciar}
En \emph{Mi trámite de licenciatura}, bajo \emph{¿Cómo quieres titularte?}, abre \emph{Sin trabajo escrito} y elige \emph{Examen General de Egreso (EGEL)}. La página de la opción explica qué necesitas, cómo será tu trámite y ofrece este manual para descargar, sin iniciar nada. Para empezar, acepta el aviso de privacidad y pulsa \boton{Iniciar mi trámite}. En EGEL, antes eliges la convocatoria.
\captura{egel_pasante_01_entrada}{«¿Cómo quieres titularte?»: abres «Sin trabajo escrito» y eliges la opción.}
\captura{egel_pasante_02_informativa}{Página de la opción, antes de iniciar: qué necesitas, cómo será el trámite y el manual para descargar.}
\captura{egel_pasante_03_iniciar}{Aceptas el aviso de privacidad y pulsas «Iniciar mi trámite».}
\subsection{Tu trámite, etapa por etapa}
Tu trámite tiene 6 etapas. Arriba siempre ves la tarjeta \emph{Tu siguiente paso} con lo que te toca ahora y una barra con las etapas: las terminadas, la que está en curso y las que vienen después.
\subsubsection*{Etapa 1. Tus datos}
Captura una sola vez tus datos; con ellos se llenan automáticamente todos los formatos del trámite.

Datos que capturas:
\begin{multicols}{2}\small
\begin{itemize}
  \item Nombre completo (con acentos, tal como aparecerá en el título)
  \item Número de cuenta
  \item Licenciatura
  \item Sistema
  \item Plan de estudios (clave y versión)
  \item Escuela incorporada a la UAEM donde cursaste la licenciatura
  \item Año de ingreso
  \item Año de egreso
  \item Teléfono
  \item Correo institucional
  \item Correo personal
  \item Ocupación actual
  \item Fecha de nacimiento
  \item Mi nombre o apellidos llevan acentos
\end{itemize}
\end{multicols}
\captura{egel_pasante_06_siguiente}{Tu siguiente paso: Tus datos}
\captura{egel_pasante_07_captura_datos}{Formulario de la etapa: capturas tus datos una sola vez.}
\subsubsection*{Etapa 2. Solicitud para presentar el EGEL}
Debes tramitar la solicitud antes de presentar el examen. Después te registras en el sitio de la Dirección de Estudios Profesionales (Pagando menos) o en la Facultad de Contaduría (convocatoria nacional).

Datos que capturas:
\begin{multicols}{2}\small
\begin{itemize}
  \item La solicitud se presenta por
\end{itemize}
\end{multicols}

La plataforma llena con tus datos los formatos de esta etapa; los firmas electrónicamente con un botón.

Documentos (PDF):
{\small\begin{longtable}{>{\raggedright\arraybackslash}p{0.47\linewidth}>{\raggedright\arraybackslash}p{0.45\linewidth}}
\toprule
\textbf{Documento} & \textbf{Observaciones}\\
\midrule
\endhead
Solicitud firmada por el pasante & La plataforma lo genera.\\
Certificado total de estudios de licenciatura & \\
Certificado de servicio social & \\
\bottomrule
\end{longtable}}
\captura{egel_pasante_08_siguiente}{Tu siguiente paso: Solicitud para presentar el EGEL}
\captura{egel_pasante_09_captura_solicitud}{Formulario de captura de la etapa.}
\captura{egel_pasante_11_generar}{La plataforma llena el formato con tus datos y lo firma por ti con un botón.}
\captura{egel_pasante_13_subir}{Subir un documento: eliges el PDF y lo envías; queda «En validación».}
\subsubsection*{Etapa 3. Registro y resultado del EGEL}
Registra tu examen en sede UAEM, preséntalo y sube tu resultado. Un resultado «Aún no satisfactorio» cuenta como evaluación aplazada.

Datos que capturas:
\begin{multicols}{2}\small
\begin{itemize}
  \item Fecha en que presentaste el EGEL
  \item Testimonio de desempeño obtenido
\end{itemize}
\end{multicols}

Documentos (PDF):
{\small\begin{longtable}{>{\raggedright\arraybackslash}p{0.47\linewidth}>{\raggedright\arraybackslash}p{0.45\linewidth}}
\toprule
\textbf{Documento} & \textbf{Observaciones}\\
\midrule
\endhead
Comprobante de registro y pago del EGEL & \\
Reporte individual de resultados y Testimonio de Desempeño Satisfactorio o Sobresaliente & Emitido por el CENEVAL. El Vo. Bo. y sello de la Subdirección los gestiona Apoyo a Titulación.\\
\bottomrule
\end{longtable}}
\captura{egel_pasante_16_siguiente}{Tu siguiente paso: Registro y resultado del EGEL}
\captura{egel_pasante_17_captura_ceneval}{Formulario de captura de la etapa.}
\subsubsection*{Etapa 4. Documentación para la toma de protesta}
Sube cada documento en PDF, legible y a color. Cuando todo esté validado, Apoyo a Titulación te citará para entregar originales y copias y fijar la fecha.

La plataforma llena con tus datos los formatos de esta etapa; los firmas electrónicamente con un botón.

Documentos (PDF):
{\small\begin{longtable}{>{\raggedright\arraybackslash}p{0.47\linewidth}>{\raggedright\arraybackslash}p{0.45\linewidth}}
\toprule
\textbf{Documento} & \textbf{Observaciones}\\
\midrule
\endhead
Certificado de estudios de bachillerato & \\
Constancia de autenticación de bachillerato & Se tramita en la escuela donde cursaste el bachillerato. Solo si tus estudios de bachillerato no fueron en instituciones dependientes o incorporadas a la UAEM.\\
Acta de nacimiento & Copia certificada vigente (del año en curso).\\
CURP & Debe mostrar la leyenda «Validada por el Registro Civil».\\
Constancia o certificado de aprobación del idioma inglés & Solo planes 98 en adelante, si en tu certificado no aparecen Inglés C1 e Inglés C2.\\
Constancia de no adeudo a la Biblioteca de la Facultad & El recibo de pago se tramita en Control Escolar de la Facultad; la constancia, en la Biblioteca.\\
Constancia de no adeudo a la Biblioteca Central & \\
Contra-recibo de la donación de libros a la Biblioteca de la Facultad & \\
Constancia de no adeudo a la Dirección de Control Escolar de la UAEM & Se solicita en el Portal de alumnos: controlescolar.uaemex.mx.\\
Recibo de no adeudo a la UAEM & Se solicita en el Portal de alumnos de Control Escolar.\\
Formato de datos personales (para el título) & La plataforma lo genera. La plataforma lo llena con tus datos; lo firmas electrónicamente.\\
Comprobante del sistema de Seguimiento de Egresados & Se obtiene en alumniversitario.uaemex.mx/1dae.php.\\
Dictamen de autorización de los HH. Consejos & La plataforma lo genera. Se solicita en la Subdirección Académica con el formato del Art. 5. Solo si excediste el plazo de 10 años para presentar la evaluación profesional (Art. 5, REP).\\
Recibo de pago: derechos de evaluación profesional (\$230) & Se solicita en Control Escolar de la Facultad.\\
Recibo de pago: expedición de título (\$1,187) & Se solicita en Control Escolar de la Facultad.\\
\bottomrule
\end{longtable}}
\captura{egel_pasante_20_siguiente}{Tu siguiente paso: Certificado de estudios de bachillerato}
\captura{egel_pasante_23_generar}{La plataforma llena el formato con tus datos y lo firma por ti con un botón.}
\subsubsection*{Etapa 5. Fecha, hora y lugar de la toma de protesta}
Apoyo a Titulación asigna el jurado y te comunica la fecha. Preséntate puntualmente.

\emph{Le toca al departamento:} Fecha y lugar de la toma de protesta. Titulación fija la fecha. En esta opción la comisión es siempre el Director (presidente) y el Subdirector Académico (secretario); el acta oficial se genera en otro sistema y se firma en papel.
\subsubsection*{Etapa 6. Resultado de la evaluación}
Al concluir, el departamento registra el resultado. Después podrás consultar cómo descargar tu título y cédula profesional.

\emph{Le toca al departamento:} Registrar resultado. 
\captura{egel_pasante_24_siguiente}{Tu siguiente paso: Resultado de la evaluación}
\subsection{Si te devuelven un documento}
Cuando Titulación devuelve un documento, arriba de tu trámite aparece \emph{Falta corregir un documento} con la observación, y tu siguiente paso te lleva a corregirlo. Sube el PDF corregido desde el mismo renglón.
\captura{egel_pasante_14_devuelto}{Un documento devuelto: arriba se muestra qué corregir y la observación de Titulación.}
\captura{egel_pasante_15_resubir}{Corriges el documento devuelto: lees la observación, eliges el PDF corregido y lo vuelves a enviar.}
\subsection{Al terminar}
Cuando se registra el resultado, tu trámite queda concluido. Titulación verifica tu expediente y lo envía a Servicios Escolares para la expedición del título. En \emph{Estado de mi trámite} ves en qué va.
\captura{egel_pasante_25_concluido}{Trámite concluido: el resultado queda registrado.}
\captura{egel_pasante_26_estado}{«Estado de mi trámite»: en qué va y qué sigue, por etapas.}
\subsection{Desde el celular}
La plataforma funciona igual en el teléfono: la tarjeta \emph{Tu siguiente paso} queda arriba y las etapas futuras se pliegan.
\capturasmovil{egel_pasante_04_movil_entrada}{Desde el celular: las opciones se abren con un toque.}{egel_pasante_05_movil_manual}{Desde el celular: el manual de la opción se descarga desde su página.}
